\documentclass[11pt]{article}

\usepackage[margin=1in]{geometry}
\usepackage{amsmath,amssymb,amsthm}
\usepackage{mathrsfs}
\usepackage{enumitem}
\usepackage[hidelinks]{hyperref}
\urlstyle{same}

\newtheorem{theorem}{Theorem}
\newtheorem{proposition}[theorem]{Proposition}
\newtheorem{corollary}[theorem]{Corollary}
\newtheorem{lemma}[theorem]{Lemma}
\theoremstyle{definition}
\newtheorem{definition}[theorem]{Definition}
\theoremstyle{remark}
\newtheorem{remark}[theorem]{Remark}

\newcommand{\F}{\mathbb{F}}
\newcommand{\Tr}{\operatorname{Tr}}
\newcommand{\Aut}{\operatorname{Aut}}
\newcommand{\nimmul}{\mathbin{\mathop{\otimes}\limits_{\mathrm{nim}}}}
\newcommand{\taglabel}[1]{\textnormal{\textbf{[#1]}}}
\newcommand{\PROVED}{\taglabel{PROVED}}
\newcommand{\STANDARD}{\taglabel{STANDARD MATH}}
\newcommand{\BOUNDARY}{\taglabel{BOUNDARY}}

\title{A Game-Native Model of the\\ Gold--Heisenberg Extension}
\author{a9lim}
\date{August 2026}

\begin{document}
\maketitle

\begin{abstract}
The extraspecial dictionary identifies quadratic maps $Q$ on an
$\F_2$-space with central extensions whose squares read $Q$ and whose
commutators read its polar form.  The remaining structural question for
the Gold trace family asked for such an extension built from game values
and game operations, without importing the diagonal $Q_a(e_i)$ or a unary
$Q_a$ oracle.  We give a basis-free answer.  The oriented half-polar term
\[
 \phi_{\lambda,a}(x,y)=\Tr(\lambda x y^{2^a})
\]
is a bilinear two-cocycle.  Its cocycle group on
$\F_2\times\F_{2^m}$ is represented by two nimber-game registers; the
product uses only disjunctive sum, Turning-Corners nim multiplication,
Frobenius, and trace.  Its squares are the scaled Gold form and its
commutators are the Gold polar form.  A signed-translation action on
finite bundles of short-game values makes the composition literal and
faithful.  We also correct the target's terminology: an unscaled Gold
form has a polar radical, so its full extension is extraspecial-type, not
extraspecial.  In the power-of-two nimber tower the radical has a
canonical central lift; quotienting by it yields the genuine extraspecial
core, with the first $(m,a)=(4,1)$ cell equal to $Q_8$.  The universal
cocycle algebra is kernel-checked in Lean~4.
\end{abstract}

\section{The target and the oriented half}

Let $K=\F_{2^m}$ be represented by its normal-play nimber values.  Fix
$a\geq0$ and $\lambda\in K$.  The scaled Gold trace form is
\[
 Q_{\lambda,a}(x)=\Tr_{K/\F_2}(\lambda x^{1+2^a}).
\]
Its polar form is
\[
 B_{\lambda,a}(x,y)=Q_{\lambda,a}(x+y)+Q_{\lambda,a}(x)
 +Q_{\lambda,a}(y).
\]
The structural extension problem is not solved by the symmetric form
$B$ alone: characteristic two keeps the diagonal independent.  The
missing datum is already present before symmetrization.

\begin{definition}[Gold half-polar cocycle]\label{def:phi}
Define
\[
 \phi_{\lambda,a}(x,y)=\Tr_{K/\F_2}(\lambda x y^{2^a}).
\]
Then
\begin{align*}
 \phi_{\lambda,a}(x,x)&=Q_{\lambda,a}(x),\\
 \phi_{\lambda,a}(x,y)+\phi_{\lambda,a}(y,x)
 &=B_{\lambda,a}(x,y).
\end{align*}
\end{definition}

Frobenius is additive and trace is linear, so $\phi$ is biadditive.  This
one observation supplies both the construction and its associativity.

\section{The game-native extension}

\begin{theorem}[Gold--Heisenberg game model]\label{thm:model}
\PROVED{}
On
\[
 \mathcal E_{\lambda,a}=\F_2\times K
\]
put
\begin{equation}\label{eq:product}
 (s,x)\diamond(t,y)=
 \bigl(s+t+\phi_{\lambda,a}(x,y),\ x+y\bigr).
\end{equation}
Let $z=(1,0)$ and $\pi(s,x)=x$.  Then:
\begin{enumerate}[label=(\roman*),leftmargin=*]
\item $\mathcal E_{\lambda,a}$ is a group, and
\[
 1\longrightarrow\langle z\rangle\longrightarrow
 \mathcal E_{\lambda,a}\xrightarrow{\ \pi\ }K^+
 \longrightarrow1
\]
is a central extension with two-point fibers.  The identity is $(0,0)$
and
\[
 (s,x)^{-1}=(s+Q_{\lambda,a}(x),x).
\]
\item Its squares and commutators are
\begin{align}
 (s,x)^2&=z^{Q_{\lambda,a}(x)},\label{eq:square}\\
 [(s,x),(t,y)]&=z^{B_{\lambda,a}(x,y)}.\label{eq:comm}
\end{align}
Thus $\mathcal E_{\lambda,a}$ is exactly the extension
$E_{Q_{\lambda,a}}$; every coordinate value $Q_{\lambda,a}(e_i)$ is
produced as the square of the lift $(0,e_i)$.
\item The product is one fixed game construction.  Represent $s$ by the
impartial value $0$ or $\ast$, and $x$ by the nimber game $\ast x$.  In
these two registers,\footnote{The first register is restricted to the
prime subfield $\{0,\ast\}$; the trace expression lands there
structurally.}
\begin{equation}\label{eq:gameproduct}
 (C,X)\diamond(D,Y)=
 \left(C\oplus D\oplus
 \Tr_{\rm game}\!\left(
 (\ast\lambda)\nimmul X\nimmul\operatorname{Frob}^{a}(Y)\right),
 X\oplus Y\right).
\end{equation}
Here $\oplus$ is disjunctive sum of impartial values,
$\nimmul$ is Conway's Turning-Corners product, Frobenius is repeated
diagonal Turning Corners, and trace is the disjunctive sum of the
Frobenius iterates.  The formula contains no diagonal table, unary
$Q_{\lambda,a}$ call, outcome test, or input-dependent branch.
\item The center is
\[
 Z(\mathcal E_{\lambda,a})=
 \{(s,r):s\in\F_2,\ r\in\operatorname{rad}B_{\lambda,a}\}.
\]
Consequently the group is extraspecial exactly when the polar form is
nondegenerate and has positive rank.  Every positive-rank nonsingular
complement $W$ of the radical gives an extraspecial subgroup
$\F_2\times W$ of order $2^{1+\dim W}$, with squaring form
$Q_{\lambda,a}|_W$.
\end{enumerate}
\end{theorem}

\begin{proof}
Biadditivity gives the normalized cocycle identity
\[
 \phi(x,y)+\phi(x+y,w)=\phi(y,w)+\phi(x,y+w).
\]
The two bracketings of three factors in~\eqref{eq:product} have the same
vector coordinate, and this identity makes their central coordinates
equal.  Thus the product is associative.  The asserted identity and
inverse follow directly because both additive coordinates have exponent
two.  The element $z$ is central of order two, $\pi$ is onto, and each
fiber is the pair $(s,x)$, $z\diamond(s,x)$.

Putting the same input in both slots gives
\[
 (s,x)^2=(\phi(x,x),0)=(Q_{\lambda,a}(x),0),
\]
which proves~\eqref{eq:square}.  Swapping the order of two factors
changes only the central coordinate, by
$\phi(x,y)+\phi(y,x)=B_{\lambda,a}(x,y)$, proving~\eqref{eq:comm}.
The standard central-extension dictionary now identifies the group with
$E_{Q_{\lambda,a}}$.

Every operation in~\eqref{eq:gameproduct} is the displayed game
realization of the corresponding field operation, so the same calculation
proves the game-register statement.  Finally, $(s,x)$ commutes with every
$(t,y)$ exactly when $B(x,y)=0$ for every $y$, which is the center
formula.  On a positive-rank nonsingular complement the center is
$\langle z\rangle$, a nonzero commutator generates it, and the finite
$2$-group identity $\Phi(G)=G^2[G,G]$ gives
$Z(G)=[G,G]=\Phi(G)=\langle z\rangle$.
\end{proof}

The fixed binary half $\phi$ is essential.  Squaring a lift necessarily
reveals $Q$ in every model of $E_Q$, so ``without a $Q$ evaluator'' cannot
coherently forbid multiplication from computing the diagonal of its own
product.  The nonvacuous condition is the one Theorem~\ref{thm:model}
meets: multiplication must not import $Q$, its basis diagonal, or an
outcome-labelled lookup table.

\section{A literal composition on game-valued contexts}

The registered product already answers the stated question.  There is
also a faithful operator model showing that its multiplication is actual
composition of game-built contexts rather than an abstract group law with
game labels attached.

Let $\mathscr G$ be the additive group of short partizan game values and
$\mathscr C=\mathscr G^K$, the finite bundle of game values indexed by
the nimber field.  Define
\begin{equation}\label{eq:contexts}
 (zF)(u)=-F(u),\qquad
 (T_xF)(u)=(-1)^{\phi_{\lambda,a}(u,x)}F(u+x).
\end{equation}
The coordinate translation uses nim addition; the sign is ordinary game
negation.

\begin{proposition}[faithful signed translations]\label{prop:contexts}
\PROVED{}
With composition ordered as $T_xT_y=T_x\circ T_y$,
\[
 T_xT_y=z^{\phi_{\lambda,a}(x,y)}T_{x+y}.
\]
Hence $(s,x)\mapsto z^sT_x$ is a faithful homomorphism
\[
 \mathcal E_{\lambda,a}\hookrightarrow\Aut(\mathscr C).
\]
\end{proposition}

\begin{proof}
Biadditivity gives
\begin{align*}
 (T_xT_yF)(u)
 &=(-1)^{\phi(u,x)+\phi(u+x,y)}F(u+x+y)\\
 &=(-1)^{\phi(u,x+y)+\phi(x,y)}F(u+x+y),
\end{align*}
which is the multiplication law.  For faithfulness, let $\delta_0$ take
the integer-game value $1$ at $0$ and $0$ elsewhere.  The support of
$z^sT_x\delta_0$ is the singleton $\{x\}$, and its value there has sign
$(-1)^{s+Q_{\lambda,a}(x)}$.  Thus the image determines both $x$ and
$s$.
\end{proof}

This is a representation of the structural multiplication, not the
forbidden Cayley move rule $g\to gn$ with fixed $n$: play still uses the
flagship paper's automorphism-equivariant squaring rule once the extension
is built.

\section{Radical discipline and the genuine extraspecial core}

The phrase ``extraspecial $E_{Q_a}$'' hid a mathematical defect.  The
full unscaled Gold polar form is always singular.  The central extension
above is canonical on the full field; extraspeciality belongs to its
nonsingular core.

\begin{proposition}[Gold radical]\label{prop:radical}
\STANDARD{}
For the unscaled form $Q_a$ on $K=\F_{2^m}$, put
$d=\gcd(2a,m)$.  Then
\[
 \operatorname{rad}B_a=\F_{2^d},\qquad
 \operatorname{rank}B_a=m-d.
\]
In particular every unscaled Gold extension has center larger than
$\langle z\rangle$; when $m=2^k$ and $a\geq1$, it is never itself
extraspecial.
\end{proposition}

\begin{proof}
Frobenius invariance and trace adjointness give, with exponents read modulo
$m$,
\[
 B_a(x,y)=\Tr\!\left(
 (x^{2^{m-a}}+x^{2^a})y\right).
\]
The trace pairing is nondegenerate.  Hence $x$ is radical exactly when
$x^{2^{m-a}}=x^{2^a}$, equivalently $x^{2^{2a}}=x$.  The fixed field is
$\F_{2^{\gcd(2a,m)}}$.
\end{proof}

For a general degree, if $Q_a$ is nonzero on the radical then no quadratic
form descends to the quotient; one must choose a nonsingular complement,
and its Arf type may depend on that choice.  The power-of-two tower used by
the project lies in the cleaner case.

\begin{theorem}[canonical power-of-two radical quotient]\label{thm:quotient}
\PROVED{}
Assume $m$ is a power of two, $a\geq1$, and
$d=\gcd(2a,m)<m$.  Put $R=\F_{2^d}$ and
\[
 H=\{(0,r):r\in R\}\subseteq\mathcal E_a.
\]
Then $Q_a|_R=0$ and $\phi_a|_{R\times R}=0$.  The subgroup $H$ is central
and elementary abelian, and
\[
 1\longrightarrow\langle\bar z\rangle\longrightarrow
 \mathcal E_a/H\longrightarrow K/R\longrightarrow1
\]
is a genuine extraspecial extension of order $2^{1+m-d}$ with the induced
nonsingular Gold quadratic form.  For $(m,a)=(4,1)$ it is $Q_8$.
\end{theorem}

\begin{proof}
The relative degree $m/d$ is even.  For $r,s\in R$, both
$rs^{2^a}$ and $r^{1+2^a}$ lie in $R$, while absolute trace from $K$
restricts on $R$ to $(m/d)$ times the absolute trace of $R$.  It therefore
vanishes, proving both restrictions.  Proposition~\ref{prop:radical}
makes $H$ central, and $\phi|_{R\times R}=0$ makes it a subgroup.  The
quotient quadratic form is well defined because $Q|_R=B(R,K)=0$; its
polar form has zero radical by construction.  Theorem~\ref{thm:model}
therefore makes the quotient extraspecial and gives its order.  At
$(4,1)$ the rank is two and the already established Gold core has Arf
invariant one, so the order-eight group is the minus cell $Q_8$.
\end{proof}

Thus the proposed $Q_8$ ordinal-sum search is unnecessary: the quaternion
cell is the first canonical radical quotient of the trace-cocycle model.
On any anisotropic plane more directly, lifts $i,j$ satisfy
$i^2=j^2=z$ and $ij=zji$, which is the standard presentation of $Q_8$;
hyperbolic planes give $D_8$.

\section{Exact boundary and formal status}

\begin{remark}[naturality boundary]
\BOUNDARY{}
For the unscaled form, the construction is basis-free and Frobenius
covariant:
\[
 \phi_a(x^{2^j},y^{2^j})=\phi_a(x,y).
\]
It is not ambient-coherent under even-degree subfield inclusions, because
absolute trace restricts to zero there.  This agrees with, rather than
evades, the flagship paper's divisibility no-go for one coefficient-faithful
datum on the full additive group of short games.  Nor can a lift section
be fixed by every extension automorphism: central twists exchange its two
lifts.  Neither stronger condition belonged to the model problem.
\end{remark}

Theorem~\ref{thm:model} satisfies the stated target literally: its
positions are pairs of game values, its multiplication is a fixed
composition of named game constructions, and its diagonal is produced by
squaring.  A stronger demand that the carrier be one untagged Conway value
under one pre-existing binary compound would be a different problem.  It
would also require a precise congruence: ordinal sum does not descend to
ordinary game values, while arbitrary quotients of literal game forms can
import any finite presentation and are therefore too permissive.

The universal algebra is kernel-checked in
\path{formal/Ogdoad/GoldExtraspecial.lean}.  For an arbitrary biadditive
$\phi:V\to V\to\F_2$, Lean constructs the group, proves cocycle
associativity, central retagging, the diagonal square identity, the
symmetrized swap/commutator defect, the exact center criterion,
noncommutativity from nonzero polar value, and order-four lifts from a
diagonal-one input.  It does not encode the concrete $u128$ nimber field,
Turning-Corners implementation, or finite-field trace; those are the
paper-level specialization and the independently tested Rust surface.

\begin{thebibliography}{9}\footnotesize
\bibitem{goldarf} a9lim, \emph{Quadratic refinements in normal play:
realization and the observation boundary}, companion flagship paper,
\path{writeups/goldarf.tex} in the ogdoad repository, 2026.
\bibitem{Gorenstein80} D.~Gorenstein, \emph{Finite Groups}, 2nd ed.,
Chelsea Publishing, New York, 1980.
\bibitem{Quillen71} D.~Quillen, \emph{The mod~2 cohomology rings of
extra-special 2-groups and the spinor groups}, Math. Ann. \textbf{194}
(1971), 197--212.
\end{thebibliography}

\end{document}
